% =====================================================================
%  Section 3 — Privacy Model and Channel-Pruning Feature Release
%
%  Drop-in LaTeX for the privacy / methods section of the WACV paper.
%  No proprietary packages: needs only amsmath, amssymb, booktabs.
%
%  Numbers in Table 1 come from:
%      phase1_distillation/scripts/privacy_accounting.py
%  Numbers in Tables 2--3 come from local DRIVE experiments
%  (drive_local_demo.py, drive_wf_threshold.py, drive_student_adapter.py
%  in phase1_distillation/).
% =====================================================================

\section{Privacy Model and Channel-Pruning Feature Release}
\label{sec:privacy}

We formalise the privacy model and the proposed mechanism in five steps.
\S\ref{ssec:threat} states the threat model and explains why it differs
from the implicit \emph{per-iteration release} model used by prior DP
feature distillation. \S\ref{ssec:accounting} gives the composition
accounting. \S\ref{ssec:wf} derives the channel-importance water-filling
allocation (Theorem~\ref{thm:wf}) and shows empirically why \emph{plain}
WF is essentially equivalent to uniform noise on trained
bottleneck-feature decoders. \S\ref{ssec:thr} introduces the
\emph{threshold} mechanism that we propose: it withholds low-importance
channels entirely, concentrates the privacy budget on a small protected
subset, and recovers most of the lost utility
(Theorem~\ref{thm:thr-optimal}). \S\ref{ssec:adapter} describes the
public-proxy adapter that absorbs the residual feature error via pure
feature matching, requiring no ground-truth supervision and remaining
DP-honest by post-processing.

% ---------------------------------------------------------------------
\subsection{Threat model: sample-once-per-image release}
\label{ssec:threat}

We consider a federated medical setting in which a data-holding site
processes each patient image \emph{once} through the teacher encoder,
adds Gaussian noise to the resulting bottleneck representation, and
releases the noisy representation. A separate party trains the student
against these fixed noisy targets. Formally, for training set
$\mathcal{D}=\{x_i\}_{i=1}^{N}$ and teacher encoder
$\phi_T : \mathcal{X}\!\to\!\mathbb{R}^{C\times H\times W}$, the
released objects are
\begin{equation}
\widetilde{z}_i \;=\; \mathrm{denorm}\!\bigl(\,\mathrm{norm}(\phi_T(x_i)) + \eta_i\,\bigr),
\qquad \eta_i \sim \mathcal{N}(0,\,\mathrm{diag}(\sigma^2)),
\label{eq:release}
\end{equation}
where $\mathrm{norm}(\cdot)$ performs per-channel $L_2$ clipping and
normalisation so that each channel of the released tensor lies in the
unit ball before noise is added.

We assume \textbf{per-patient privacy with at most one image per
patient} (DRIVE: $k\!=\!1$; BraTS: $k\!=\!1$; for multi-image-per-patient
sets such as MIMIC-CXR, $k$-wise parallel composition is applied per
patient). Both images and ground-truth labels are treated as private;
the student does not access labels (\S\ref{ssec:adapter}), so the
released features are the only object the privacy mechanism must
protect.

\textbf{Why this differs from prior work.}
Existing implementations of DP feature distillation re-sample noise
\emph{at every training iteration} of the student
\cite{lan2024gkd,Anonymous_anonymous}. Under this implicit
\emph{per-iteration} model, each training image's features are
released $Q = TB/N$ times across $T$ iterations of batch size $B$;
adaptive composition over $Q$ releases blows the user-level $\varepsilon$
up by orders of magnitude (cf.\ Table~\ref{tab:privacy}). The
sample-once model in \eqref{eq:release} eliminates this composition
entirely: each user contributes exactly one release.

% ---------------------------------------------------------------------
\subsection{Composition and budget accounting}
\label{ssec:accounting}

We use zero-Concentrated Differential Privacy (zCDP)
\cite{bun2016concentrated} as the working unit. For a Gaussian
mechanism on a $C$-dimensional vector with per-coordinate sensitivity
$\Delta_c$ and per-coordinate noise scale $\sigma_c$, the per-release
zCDP cost is
\begin{equation}
\rho_{\mathrm{rel}}
 \;=\; \sum_{c=1}^{C}\frac{\Delta_c^2}{2\,\sigma_c^2}.
\label{eq:rho-rel}
\end{equation}
After per-channel $L_2$ normalisation each channel has norm at most
$1$, so the replace-one $L_2$ sensitivity is the data-independent
constant $\Delta_c = 2/K$ for all $c$, where $K\geq 1$ is the
PATE-style teacher ensemble size (we use $K=1$ in our main results;
$K=10$ is reported as a sensitivity study).

\paragraph{Conversion to $(\varepsilon,\delta)$-DP.}
By \cite[Prop.~1.3]{bun2016concentrated}, a $\rho$-zCDP mechanism is
$(\rho + 2\sqrt{\rho\log(1/\delta)},\,\delta)$-DP for any $\delta>0$.
Conversely, the inverse map
\begin{equation}
\rho(\varepsilon,\delta)
 \;=\;\Bigl(\sqrt{\log(1/\delta)+\varepsilon}-\sqrt{\log(1/\delta)}\,\Bigr)^{2}
\label{eq:eps-to-rho}
\end{equation}
gives the maximum zCDP that fits inside a target
$(\varepsilon,\delta)$-DP budget.

\paragraph{Composition.}
The two threat models differ only in the number of times each user's
data participates in a release.
\begin{itemize}
\item \textbf{Sample-once per image} (ours). Each user contributes a
single release. The user-level cost is
$\rho_{\mathrm{user}} = \rho_{\mathrm{rel}}$, giving
$\varepsilon_{\mathrm{user}} = \varepsilon_{\mathrm{rel}}$. Because
the per-channel importance and caps are computed on a held-out public
dataset (\S\ref{ssec:adapter}), no additional release accounting is
required.
\item \textbf{Per-iteration release} (prior work's implicit model).
Each user participates in $Q = TB/N$ releases of fresh noise on fresh
teacher queries. Naive zCDP composition gives
$\rho_{\mathrm{user}} = Q\,\rho_{\mathrm{rel}}$, mapped back to
$(\varepsilon_{\mathrm{user}}, \delta)$ via
\cite[Prop.~1.3]{bun2016concentrated}.
\end{itemize}

\paragraph{Reported vs.\ true privacy.}
Table~\ref{tab:privacy} contrasts the user-level $\varepsilon$ implied
by each threat model for the configurations used throughout this paper
($T{=}40{,}000$, $B{=}4$, $N{=}20$, $\delta{=}10^{-5}$). The
per-iteration model inflates the privacy cost by 400--1800$\times$
relative to what is reported by the noise calibration; the sample-once
model in \eqref{eq:release} faithfully reports user-level privacy by
construction.

\begin{table}[t]
\centering
\small
\caption{User-level $\varepsilon$ implied by each threat model for a
range of per-release $\varepsilon$ values. The per-iteration model
(prior work's implicit model) yields privacy guarantees that are
hundreds to thousands of times weaker than reported; the sample-once
model (ours) reports the actual user-level guarantee.
Parameters: $T{=}40{,}000$, $B{=}4$, $N{=}20$, $\delta{=}10^{-5}$.}
\label{tab:privacy}
\begin{tabular}{c c r r}
\toprule
per-release $\varepsilon$ & threat model & releases/user & user-level $\varepsilon$ \\
\midrule
\multirow{2}{*}{$2$}  & per-iteration & $8000$ & $812.1$ \\
                      & sample-once   & $1$    & $2.0$   \\
\addlinespace[2pt]
\multirow{2}{*}{$4$}  & per-iteration & $8000$ & $2712.4$ \\
                      & sample-once   & $1$    & $4.0$    \\
\addlinespace[2pt]
\multirow{2}{*}{$8$}  & per-iteration & $8000$ & $9014.8$ \\
                      & sample-once   & $1$    & $8.0$    \\
\addlinespace[2pt]
\multirow{2}{*}{$16$} & per-iteration & $8000$ & $28569.6$ \\
                      & sample-once   & $1$    & $16.0$    \\
\bottomrule
\end{tabular}
\end{table}

% ---------------------------------------------------------------------
\subsection{Channel-importance water-filling and its empirical limit}
\label{ssec:wf}

Equation \eqref{eq:rho-rel} shows that a fixed per-release budget
$\rho_{\mathrm{rel}}$ can be spent across channels in many different
ways. The canonical choice minimises the
\emph{importance-weighted reconstruction error}
\begin{equation}
\min_{\sigma_1,\ldots,\sigma_C > 0}
\sum_{c=1}^C s_c\,\sigma_c^{2}
\quad\text{s.t.}\quad
\sum_{c=1}^C\frac{\Delta_c^{2}}{2\,\sigma_c^{2}} \le \rho_{\mathrm{rel}},
\label{eq:opt}
\end{equation}
where $s_c \ge 0$ is a per-channel importance score (we use
$s_c = \mathrm{mean}_{i,j} |\partial\mathcal{L}_{\mathrm{task}}/\partial \phi_{T,c}^{i,j}|$,
the GKD-style gradient magnitude \cite{lan2024gkd}, computed on the
public proxy of \S\ref{ssec:adapter}). The Lagrangian solution is
closed-form.

\begin{theorem}[Channel-WF optimal allocation]
\label{thm:wf}
For positive importance scores $\{s_c\}$, positive sensitivities
$\{\Delta_c\}$ and budget $\rho_{\mathrm{rel}}>0$, the unique
minimiser of \eqref{eq:opt} is
\begin{equation}
\sigma^{\star}_c \;=\; \kappa \,\sqrt{\Delta_c}\, s_c^{-1/4},
\qquad
\kappa \;=\; \sqrt{\frac{1}{2\rho_{\mathrm{rel}}}
                          \sum_{c=1}^{C}\Delta_c\,\sqrt{s_c}}.
\label{eq:wf}
\end{equation}
\end{theorem}

\begin{proof}[Proof sketch]
Lagrangian stationarity in $u_c = 1/\sigma_c^2$ gives
$u_c \propto \sqrt{s_c}/\Delta_c$; substituting into the active
constraint solves for the multiplier. The minimum scales as
$\bigl(\sum_c \Delta_c \sqrt{s_c}\bigr)^2 / (2\rho_{\mathrm{rel}})$,
which is strictly below the uniform value
$\bigl(\sum_c \Delta_c^2\bigr)\bigl(\sum_c s_c\bigr) / (2\rho_{\mathrm{rel}})$
whenever the $s_c$ are not all equal (Cauchy--Schwarz).
The full proof appears in Appendix~\ref{app:proof}.
\end{proof}

\paragraph{Two empirical limits of plain WF.}
We probe the mechanism on the bottleneck of a U-Net teacher trained on
the DRIVE retinal-vessel dataset
(\S\ref{sec:experiments}). Two phenomena consistently violate the
intuition motivating \eqref{eq:wf}:
\begin{enumerate}
\item\textbf{The exponent $\tfrac14$ in \eqref{eq:wf} compresses the
allocation aggressively.} A per-channel importance ratio
$s_{\max}/s_{\min} = R$ yields a noise-scale ratio of only $R^{1/4}$.
On real DRIVE we measure $R\!\approx\!2$, giving $R^{1/4}\!\approx\!1.19$
\textemdash{} plain WF is operationally indistinguishable from uniform.
A controlled ablation (Fig.~\ref{fig:imp-sweep})
that artificially boosts $R$ to $10^{5}$ ($R^{1/4} \!=\!10$) shows
plain WF actually \emph{loses} to uniform at this extreme by up to
$0.01$ Dice points.

\item\textbf{Trained decoders use channels jointly.} The gradient
weight $s_c$ measures the local sensitivity of the task loss to a
clean channel, not the consequence of polluting that channel with
Gaussian noise. Pollution on a low-$s_c$ channel still propagates
through convolutional mixing in the decoder and degrades the
prediction. Concentrating noise on ``unimportant'' channels (the
behaviour of \eqref{eq:wf}) is therefore not free.
\end{enumerate}
Together these two effects explain a robust empirical finding:
\textit{on the bottleneck of a trained segmentation network, plain WF
matches uniform Gaussian to within Dice~$\pm\,0.003$ at every
$\varepsilon$ tested}, regardless of the importance ratio.
\S\ref{ssec:thr} resolves this by changing what plain WF spends the
budget on.

% ---------------------------------------------------------------------
\subsection{Threshold mechanism: channel pruning under DP}
\label{ssec:thr}

The empirical failure of plain WF is structural, not numerical:
spending budget on channels whose pollution still confuses the
decoder cannot be optimal. Our proposed mechanism removes those
channels from the release entirely.

\paragraph{Mechanism.}
Given a public-proxy importance vector $\widehat{s}\in\mathbb{R}^C$
and a keep-fraction $k\in(0,1]$, define the active set
$\mathcal{A}\subseteq\{1,\dots,C\}$ as the top
$\lfloor kC \rfloor$ channels by $\widehat{s}_c$. Allocate the entire
budget $\rho_{\mathrm{rel}}$ on $\mathcal{A}$, set
$\sigma_c=\infty$ on the complement, and replace the off-$\mathcal{A}$
channels with the data-independent constant $0$ in the normalised
representation:
\begin{equation}
\widetilde{z}_{i,c} \;=\;
\begin{cases}
 \widehat{c}_c\bigl(\,\mathrm{norm}(\phi_T(x_i))_{c} + \eta_{i,c}\,\bigr),
 & c\in\mathcal{A},\\[2pt]
 0, & c\notin\mathcal{A},
\end{cases}
\qquad
\eta_{i,c}\sim\mathcal{N}(0,\sigma_c^{2}\mathbf{I}).
\label{eq:thr-release}
\end{equation}
The per-channel zCDP cost \eqref{eq:rho-rel} sums only over
$\mathcal{A}$ (inactive channels contribute $\Delta_c^2/(2\sigma_c^2)=0$),
so the full budget is concentrated on $|\mathcal{A}|$ channels rather
than spread across $C$.

\paragraph{An empirical surprise: WF and uniform on $\mathcal{A}$
coincide.}
On the active set we have two natural choices: solve \eqref{eq:opt}
restricted to $\mathcal{A}$ (a WF solution), or assign the uniform
$\sigma^{\mathrm{unif}}_{\mathcal{A}} =
 \sqrt{(\sum_{c\in\mathcal{A}}\Delta_c^2)/(2\rho_{\mathrm{rel}})}$.
On real DRIVE these two choices give identical Dice to four decimal
places at every $\varepsilon\!\in\!\{4,8,16,32\}$ and every keep-fraction
$k\!\in\!\{0.1,0.25\}$ that we test
(Table~\ref{tab:thr-vs-unif-on-active}; Fig.~\ref{fig:k1-and-uniform-thr}b).
The intuitive explanation is that
the importance ratio \emph{within the active set} is even flatter
than across all $C$, so the $\tfrac14$-power closed form is
essentially constant on $\mathcal{A}$.

This justifies using the simpler uniform allocation on $\mathcal{A}$.
We refer to the mechanism as \emph{channel-pruning feature release}
(CPFR) when emphasising the operational form
(uniform Gaussian on $\mathcal{A}$, zero off $\mathcal{A}$), and as
\emph{WF$+$threshold} when emphasising the connection to
Theorem~\ref{thm:wf}.

\paragraph{Optimality under sparse importance.}
The threshold mechanism is provably optimal in a clean limit: when
the importance vector has at most $k$ non-zero coordinates, the
solution of the augmented optimisation (allowing $\sigma_c = \infty$)
sets $\sigma_c = \infty$ on every zero-importance coordinate.

\begin{theorem}[Threshold optimality under sparse importance]
\label{thm:thr-optimal}
Let $\mathcal{S}=\{c: s_c>0\}$ and $C^\star=|\mathcal{S}|$. Consider
the augmented problem
\begin{equation*}
\min_{\sigma_1,\dots,\sigma_C\in(0,\infty]}
 \sum_{c=1}^C s_c\,\sigma_c^{2}
 \quad\text{s.t.}\quad
 \sum_{c=1}^C \frac{\Delta_c^2}{2\sigma_c^2} \le \rho_{\mathrm{rel}},
\end{equation*}
with the conventions $s_c\sigma_c^2 = 0$ when $s_c=0$ and
$\Delta_c^2/\sigma_c^2 = 0$ when $\sigma_c=\infty$. Then every
optimal solution satisfies $\sigma_c=\infty$ for $c\notin\mathcal{S}$,
and on $\mathcal{S}$ the optimum is given by Theorem~\ref{thm:wf}
restricted to $\mathcal{S}$ with budget $\rho_{\mathrm{rel}}$. The
optimal objective is
\begin{equation}
J^\star_{\mathrm{thr}}
 \;=\;
 \frac{1}{2\rho_{\mathrm{rel}}}
 \Bigl(\sum_{c\in\mathcal{S}} \Delta_c \sqrt{s_c}\Bigr)^2,
\label{eq:thr-obj}
\end{equation}
strictly less than $J^\star_{\mathrm{WF}}$ over all $C$ channels by a
factor of $(\sum_{c\in\mathcal{S}}\Delta_c\sqrt{s_c})^2 /
(\sum_{c=1}^C\Delta_c\sqrt{s_c})^2 < 1$.
\end{theorem}

\begin{proof}[Proof sketch]
Inactive coordinates contribute $0$ to the objective and need
$0$ budget at $\sigma_c=\infty$; any finite $\sigma_c$ on an inactive
coordinate spends budget without benefit. Restrict to $\mathcal{S}$
and apply Theorem~\ref{thm:wf}. The objective ratio follows from
\eqref{eq:wf-obj}. Full proof in Appendix~\ref{app:thr-proof}.
\end{proof}

\paragraph{From sparse importance to thresholding.}
Real importance vectors are not exactly sparse: $s_c > 0$ for every
$c$. The threshold mechanism approximates the sparse case by
zeroing the smallest $(1-k)C$ entries before solving
Theorem~\ref{thm:wf}. The empirical equivalence of WF and uniform on
$\mathcal{A}$ (above) further simplifies the mechanism to its
practical form \eqref{eq:thr-release} with $\sigma_c\equiv\sigma^{\mathrm{unif}}_{\mathcal{A}}$.

\paragraph{Empirical lift over plain WF.}
On real DRIVE (Table~\ref{tab:main-mech}, Fig.~\ref{fig:tradeoff}),
CPFR at $k=0.1$ improves vessel Dice over plain WF by $+0.052$ at
$\varepsilon=8$, $+0.094$ at $\varepsilon=16$, and $+0.168$ at
$\varepsilon=32$ (K=10) \textemdash{} the only mechanism in our sweep
that meaningfully reduces the privacy--utility gap on bottleneck
feature release with a trained decoder.

% ---------------------------------------------------------------------
\subsection{Public-proxy adapter (DP-honest post-processing)}
\label{ssec:adapter}

A residual gap to clean Dice remains because the active-set noise is
not zero. We close part of it with a small adapter, trained
\emph{without any ground-truth supervision and without any access to
private DRIVE data}, that maps a noisy thresholded release back to
the clean bottleneck space.

\paragraph{Adapter and training.}
Let $\phi_T^{(\mathrm{pub})}$ denote applying the (frozen) teacher
encoder to a public retinal proxy dataset
$\mathcal{D}_{\mathrm{pub}}$ (HRF in our experiments). For each
public image $x'$ we form the pair
$\bigl(\widetilde{z}'(\sigma), \phi_T(x')\bigr)$ where
$\widetilde{z}'(\sigma)$ applies \eqref{eq:thr-release} to $x'$ with
the same $\sigma$ used for the private release. A small residual
3-layer convolutional adapter $h_\theta$ is trained by
\begin{equation}
\min_\theta\;\mathbb{E}_{x'\sim\mathcal{D}_{\mathrm{pub}}}\;
 \bigl\|\,h_\theta\!\bigl(\widetilde{z}'(\sigma)\bigr) - \phi_T(x')\,\bigr\|_2^{2}.
\label{eq:adapter}
\end{equation}
At inference on the private set: the released $\widetilde{z}_i$ is
passed through $h_\theta$ and then through the frozen teacher decoder.

\paragraph{Privacy.}
$h_\theta$ depends only on (i) the (frozen) teacher's response on
public data and (ii) the public mechanism $\sigma$. Composition with
the private release \eqref{eq:thr-release} is therefore pure
post-processing, contributing zero additional privacy cost. The
adapter \emph{never sees ground-truth vessel masks of any kind}.

\paragraph{Empirical lift over no-adapter probe.}
On real DRIVE (Fig.~\ref{fig:adapter}), the adapter recovers an
additional $+0.081$ Dice at $\varepsilon=8$, $+0.065$ at
$\varepsilon=16$, and $+0.050$ at $\varepsilon=32$, closing
approximately $19\%$ of the residual gap to clean Dice at every
operating point. Stacking adapter on top of CPFR $k=0.1$ yields the
full pipeline result of $0.39$ Dice at $\varepsilon=32$, up from $0.17$
under uniform Gaussian release \textemdash{} a $2.3\times$ utility
improvement at fixed privacy.

% ---------------------------------------------------------------------
\subsection{Implementation}
\label{ssec:impl}

The sample-once threat model and the threshold release of
\eqref{eq:thr-release} are realised by a single precompute pass
before student-adapter training: for every training image, the
encoder is run once, the bottleneck is normalised, the active set is
selected by public-proxy importance, $\sigma$ is fixed once per job
(uniform on $\mathcal{A}$), the noise is sampled, off-$\mathcal{A}$
channels are zeroed, and the result is denormalised and cached. The
adapter \eqref{eq:adapter} is trained on the public proxy in a
separate pass and used as a frozen post-processor at inference. The
inner training loop sees no fresh teacher queries, no fresh noise,
and no ground-truth masks of the private set. Pseudocode appears in
Appendix~\ref{app:pseudocode}.
